\usepackage[T1]{fontenc}
\usepackage[utf8]{inputenc}
\usepackage[a4paper, top=2.54cm, bottom=2.54cm, left=2.54cm, right=2.54cm]{geometry}
\usepackage{mathptmx}
\usepackage{microtype}
\usepackage{amssymb}
\usepackage{parskip}
\usepackage{enumitem}
\usepackage{amsmath}
\usepackage{url}

\setlength {\marginparwidth}{2cm}

\usepackage{todonotes}
\usepackage[normalem]{ulem}
\usepackage{import}

% Title and components
% Centred, normal-weight section headings (no bold, no number).
% This also applies to \tableofcontents and \listoffigures headings,
% since LaTeX typesets them via \section*{\contentsname} etc.

% Sections
\usepackage{titlesec}
\titleformat{\section}[block]
  {\normalfont\centering} % Centered, large, bold font
  {\thesection.}                       % Adds section number followed by a dot (e.g., 1.)
  {1em}                                % Space between number and title text
  {\MakeUppercase}                     % Automatically forces title text into ALL CAPS
\titlespacing*{\section}
  {0pt}                                % Left margin offset
  {5ex plus 1ex minus .2ex}          % Vertical space ABOVE (increase 2.5ex to 3ex for more space)
  {1.5ex plus .2ex}                    % Vertical space BELOW

% Paragraph text
\setlength{\parindent}{20pt}
\urlstyle{same}

% Captions
\usepackage{caption}
\renewcommand{\thetable}{\Roman{table}}
\DeclareCaptionLabelFormat{uppercase_table}{TABLE~#2}
\DeclareCaptionFormat{multiline_title}{#1\par#3}
\captionsetup[table]{
    labelformat=uppercase_table, % Uses "TABLE" + Roman numeral
    labelsep=none,               % Removes colon/period after label
    format=multiline_title,      % Puts title on line below label
    justification=centering      % Centers both lines
}

% Model section
\usepackage{tabularray}
\usepackage{float}
\usepackage{graphicx}
\usepackage{codehigh}
\usepackage[normalem]{ulem}
\UseTblrLibrary{booktabs}
\UseTblrLibrary{siunitx}
\newcommand{\tinytableTabularrayUnderline}[1]{\underline{#1}}
\newcommand{\tinytableTabularrayStrikeout}[1]{\sout{#1}}
\NewTableCommand{\tinytableDefineColor}[3]{\definecolor{#1}{#2}{#3}}

% Bibliography, reference
\usepackage[round,authoryear]{natbib}
\bibliographystyle{plainnat}
\setcounter{secnumdepth}{0}
\pagestyle{plain}


% TOC/LOF heading names
\renewcommand{\contentsname}{Table of Contents}

% Normal-weight (not bold) entries in the TOC
\makeatletter
\renewcommand*\l@section[2]{\@dottedtocline{1}{1.5em}{2.3em}{#1}{#2}}
\makeatother

\newcommand{\lblock}[2][c]{%
  \begin{tabular}[#1]{@{}l@{}}#2\end{tabular}}
\newcommand{\cblock}[2][c]{%
  \begin{tabular}[#1]{@{}c@{}}#2\end{tabular}}